\begin{abcliste}
\abc
Auf die Rakete wirken der Schub $\ssc{F}{S}$ nach oben und die Gewichtskraft nach unten; die Differenz beschleunigt sie:
\begin{align}
 \ssc{F}{S} - m\,g &= m\,a \\
 a &= \aF \\
   &= \frac{\FT}{\mR}-\g \\
   &= \a \approx \result{\aP{0}{2}}
\end{align}

\abc
Die Astronautin hat dieselbe Beschleunigung: Ihr Sitz drückt nach oben, ihre Gewichtskraft zieht nach unten.
\begin{align}
 \ssc{F}{N} - \ssc{m}{A}\,g &= \ssc{m}{A}\,a \\
 \ssc{F}{N} &= \FNF \\
   &= \mAs\cdot(\g+\a) \\
   &= \FN \approx \result{\FNP{0}{2}}
\end{align}
Das ist das \nGP{0}{2}-Fache ihrer Gewichtskraft: Sie spürt $\nGP{0}{2}\,g$.
\end{abcliste}
